\scanpage{54}
% archive foliation: 53
un Pb. est celui où $k(\eta)=L$ est de car.p (cas d'égales caractéristiques). On peut alors penser à attaquer le Pb via la théorie de Dieudonné (supposant résolu le Pb 1) et via le

\uline{Problème} n° ... Soit S un schéma noethérien normal \deleted{de car} intègre de car. p, de point générique $\eta$. Est-il vrai que le foncteur de restriction
\[
\mathrm{Cris\ loc\ lib}\ (S)\longrightarrow\mathrm{Cris\ loc\ lib}\ (\eta)
\]
(ou du moins le foncteur analogue pour F-cristaux) est pleinement fidèle ?

On est ramené facilement au cas où S est le spectre d'un anneau de la forme $\unsure{k[[t]]|k[(t)]}$, avec k alg. clos, auquel cas l'énoncé s'explicite en termes de Modules avec connexions formelles .... Il semble que si la réponse au Pb précédent est affirmative, \deleted{ainsi que} et si on sait faire la théorie de \unsure{Dieudonnén|Dieudonné} sur un anneau \uline{parfait} (avec \uline{équivalence} de catégories alors on aura démonstration du th. de Tate en égales car.

